% Context from chapters/denoising.tex; for mathematical reference, not compilation.
evels. Numerical acquisition parameters are not available for these legacy illustrations.  
}{fig-poisson-denoising}

  
\paragraph{Variance stabilization.}

Applying a thresholding estimator
\eq{
	\Dd(f) = \sum_m S_T^q(\dotp{f}{\psi_m}) \psi_m
}
directly to $f$ can perform poorly because a single threshold $T$ cannot adapt to spatially varying noise. A variance-stabilizing transform $\phi : [0,+\infty) \rightarrow \RR$, applied pixelwise, maps the counts to $\phi(f)$. An additive Gaussian white-noise model can then be a useful approximation:
\eql{\label{eq-var-stab}
	\phi(f) \approx \phi(f_0) + w
}
where $w_n\sim\Nn(0,\sigma^2)$ are independent centered Gaussians with constant variance $\si^2$.

The model~\eqref{eq-var-stab} is approximate, and its accuracy depends on the intensity range of $f_{0,n}$.
Two common variance-stabilizing transforms for Poisson noise are the Anscombe transform
\eq{
	\phi(x) = 2\sqrt{x+3/8}
}
and the Freeman--Tukey transform
\eq{
	\phi(x) = \sqrt{x+1}+\sqrt{x}.
}
Figure \ref{fig-variance-stabilization} compares the variance of $\phi(f)$ after these transformations.

\myfigure{
\image{denoising}{.5}{variance-stabilization-corrected}
}{ 
	 Exact Poisson variances of the Anscombe, Freeman--Tukey, and $2\sqrt{x}$ transforms as functions of the mean intensity.  
}{fig-variance-stabilization}

A variance-stabilized denoiser is defined as
\eq{
	\Dd^{\text{stab},q}(f) =  \phi^{-1}( \sum_m S_T^q(\dotp{\phi(f)}{\psi_m}) \psi_m )
}
where $\phi^{-1}$ is applied after projecting the denoised values onto its domain. The direct inverse generally introduces bias, especially at low counts; bias-corrected inverse transforms can improve the result.

At the moderate count levels in Figure \ref{fig-poisson-wav}, variance stabilization improves the denoising result.

\myfigure{
\tabtrois{
\image{denoising}{.32}{poisson-wav-noisy}&
\image{denoising}{.32}{poisson-wav-unstab}&
\image{denoising}{.32}{poisson-wav-stab}
}
}{ 
	Left: noisy image, center: denoising without variance stabilization, right: denoising after variance stabilization.   
}{fig-poisson-wav}

                                                     
\subsection{Multiplicative Noise}

  
\paragraph{Multiplicative image formation.}

A multiplicative noise model assumes that
\eq{
	f_n = f_{0,n} w_n
}
where $w$ is a random multiplier, observed through one realization, with $\mathbb E(w_n)=1$.
Once again, the noise level depends on the pixel value
\eq{
	\mathbb E(f_n) = f_{0,n}
	\qandq
	\text{Var}(f_n) = f_{0,n}^2 \si^2 
	\qwhereq 
	\sigma^2=\operatorname{Var}(w_n).
}
In SAR imaging, a common model forms $f$ by averaging $K$ independent observations, called looks:
\eq{
	\foralls 1\leq s\leq K, \quad f^{(s)}_n = f_{0,n} w^{(s)}_n + r^{(s)}_n
}
where $r^{(s)}$ is additive Gaussian white noise and the multiplier $w^{(s)}_n$ has an exponential distribution:
\eq{
	p_{w^{(s)}_n}(x)=e^{-x}\mathbf1_{x>0}.
}
Averaging $K$ independent looks divides the additive-noise variance by $K$. When the remaining additive component is negligible, the average is approximately
\eq{
	f_n = \frac{1}{K} \sum_{s=1}^K f^{(s)}_n \approx f_{0,n} w_n
}
where the multiplicative factor is distributed according to a Gamma distribution
\eq{
	w_n\sim\operatorname{Gamma}(\text{shape}=K,\text{rate}=K),\qquad p_{w_n}(x)=\frac{K^K}{\Gamma(K)}x^{K-1}e^{-Kx}\mathbf1_{x>0}.
}
Its mean is one and its variance is $1/K$, so increasing $K$ reduces the relative noise level.

\myfigure{
\tabquatre{
\image{denoising}{.23}{multnoise-1}&
\image{denoising}{.23}{multnoise-2}&
\image{denoising}{.23}{multnoise-3}&
\image{denoising}{.23}{multnoise-4}
}
}{ 
	 Images with multiplicative noise at different values of $\si$.  
}{fig-multnoise}

Figure \ref{fig-multnoise} illustrates how the noise changes with the number $K=1/\si^2$ of averaged looks.

A simple variance stabilization transform is
\eq{
	\phi(x)=\log(x) - c
}
where
\eq{
	c = \text{E}(\log(w)) = \psi(K) - \log(K)
	\qwhereq
	\psi(x) = \Ga'(x)/\Ga(x)
}
Here $\Ga$ is the Gamma function and $\psi$ is its logarithmic derivative, the digamma function. The transformed observation satisfies
\eq{
	\phi(f)_n=\log f_{0,n}+z_n,
}
for strictly positive intensities, where $z_n=\log(w_n)-c$ is centered additive noise with variance $\psi^{(1)}(K)$; $\psi^{(1)}$ is the derivative of the digamma function. Thus the transformed target is $\log f_0$, not $\phi(f_0)$. Exponentiating an estimate of $\log f_0$ introduces a further statistical bias and may require correction.

\myfigure{
\image{denoising}{.45}{multnoise-hist-unstab}
\image{denoising}{.45}{multnoise-hist-stab}
}{ 
	 Histogram of multiplicative noise before (left) and after (right) stabilization.   
}{fig-multnoise-hist-unstab}

Figure \ref{fig-multnoise-hist-unstab} shows the effect of this variance stabilization on the distributions of $w$ and $z$.

Figure \ref{fig-multnoise-wav} compares the two approaches at moderate noise levels $\si$, where variance stabilization improves the displayed result.

\myfigure{
\tabtrois{
\image{denoising}{.32}{multnoise-wav-noisy}&
\image{denoising}{.32}{multnoise-wav-stab}&
\image{denoising}{.32}{multnoise-wav-unstab}
}
}{ 
	Left: noisy image, center: denoising after variance stabilization, right: denoising without variance stabilization.  
}{fig-multnoise-wav}


